\documentclass[sigconf, nonacm]{acmart}

%% ── Suppress all ACM boilerplate ────────────────────────────────────────────
\setcopyright{none}
\settopmatter{printacmref=false}
\renewcommand\footnotetextcopyrightpermission[1]{}
\pagestyle{plain}

%% ── Packages ────────────────────────────────────────────────────────────────
\usepackage{booktabs}
\usepackage{amsmath}
\usepackage{microtype}
\usepackage{graphicx}
\usepackage{multirow}
\usepackage{stfloats}
\usepackage{float}

%% ── Document metadata ───────────────────────────────────────────────────────
\title{Market Reactions to Natural Disasters: An Event Study of\\
       Cyclone Ditwah's Impact on the Colombo Stock Exchange}

\author{
Jayawardana K.P.T.,
Wickramaratne W.P.G.A.,
Wijesinghe D.S.,
Fernando W.T.D.,
Dinapura H.H.S.
}

\affiliation{%
  \institution{Department of Computer Science \& Engineering, University of Moratuwa}
  \city{Moratuwa}
  \country{Sri Lanka}
}

\email{{tharupahanj.23, anupamaw.23, dehanw.23, thevinduf.23, senindud.23}@cse.mrt.ac.lk}

\begin{document}

\begin{abstract}
%% FIX: removed the sentence claiming significance testing is in a companion
%% paper — it is now fully reported here.
Natural disasters inflict severe macroeconomic shocks that ripple through
financial markets, yet the dynamics of such shocks in emerging-economy
exchanges remain understudied. Cyclone Ditwah made landfall in eastern Sri
Lanka on 28~November~2025, causing an estimated USD~4.1~billion in direct
damage (approximately 4\% of GDP) and triggering the Colombo Stock
Exchange's (CSE) worst weekly decline since November~2022. This paper
employs a classical event-study framework on a self-collected dataset of
70,413 daily OHLCV records for 268 CSE-listed stocks spanning
February~2025 to February~2026. Market-model parameters ($\hat{\alpha}$,
$\hat{\beta}$) are estimated for 237 qualifying stocks. Sector-level
cumulative abnormal returns (CARs) are computed across 20 standard CSE
sectors. Transportation and Real Estate recorded the largest full-window
CARs at $-10.03\%$ and $-9.23\%$ respectively, while Commercial Services
showed resilience with a near-zero full-window CAR of $+0.27\%$.
Parametric cross-sectional $t$-tests and non-parametric Wilcoxon
signed-rank tests identify three sectors with statistically significant
negative median CARs: Insurance ($p = 0.009$), Real Estate ($p = 0.017$),
and Diversified Financials ($p = 0.034$), while the market-wide shock is
highly significant ($p < 0.001$). An analysis of market microstructure
reveals a two-phase response: a sharp intraday volatility spike at landfall
($t = 1.84$, $p = 0.076$) followed by a persistent liquidity freeze, with
total daily trading volume falling approximately 35\% post-event
($t = 3.76$, $p < 0.001$).
\end{abstract}

\keywords{Event Study, Natural Disaster, Colombo Stock Exchange, Cyclone
Ditwah, Abnormal Returns, Sector Heterogeneity, Emerging Markets,
Disaster Economics}

\maketitle

%%=============================================================================
\section{Introduction}
%%=============================================================================

Sri Lanka's economy, still recovering under an IMF Extended Fund Facility
arrangement following the 2022 sovereign debt crisis \cite{imf2025},
suffered a major setback when Cyclone Ditwah struck the eastern coastline
on 28~November~2025. The World Bank estimated direct physical damage at
USD~4.1~billion (approximately 4\% of GDP \cite{worldbank2025}) affecting
58,340~hectares of paddy, damaging 247~km of roads and 35~bridges, and
displacing over 209,000~people across all 25~districts
\cite{who2025, ocha2025}.

Financial markets responded sharply. The CSE All-Share Price Index (ASPI)
fell 3.04\% on the first trading session after landfall, extending to a
weekly loss of 5.35\%—the worst since mid-November~2022
\cite{brecorder2025}. Net foreign outflows from the exchange accelerated
during December~2025.

Despite these observable effects, no peer-reviewed study has systematically
quantified the impact of Cyclone Ditwah on individual stocks or market
sectors. This gap matters for three stakeholder groups: investors who need
to price disaster risk into portfolios; policymakers who must design market
stabilisation interventions; and researchers who require empirical evidence
on emerging-market resilience. This paper addresses this gap using a
rigorous event-study framework \cite{mackinlay1997} applied to a
purpose-built, high-resolution CSE dataset.

This study reports on three research questions:
\begin{enumerate}
  \item \textbf{RQ1:} Did Cyclone Ditwah affect returns in the Colombo
        Stock Exchange?
  \item \textbf{RQ2:} How quickly did the market react to the cyclone?
  \item \textbf{RQ3:} Was there any predictive price reaction before the
        event?
\end{enumerate}

In addition to these primary research questions, the paper examines which
CSE sectors were most negatively affected, which were resilient or
positively impacted, whether sectors reacted at different speeds, and how
the cyclone altered broader market microstructure.

%%=============================================================================
\section{Dataset and Data Preparation}
%%=============================================================================

\subsection{Primary Dataset}

Daily OHLCV data for all CSE-listed stocks were collected via a custom
Python web-scraping pipeline querying the exchange's public data interface.
The raw dataset comprised 71,044 rows. After cleaning (detailed in
Section~2.2), the working dataset contains 70,029 rows.
Table~\ref{tab:dataset} summarises the key characteristics.

\begin{table}[h]
  \centering
  \caption{Primary dataset characteristics}
  \label{tab:dataset}
  \small
  \begin{tabular}{@{}ll@{}}
    \toprule
    \textbf{Attribute} & \textbf{Value} \\
    \midrule
    Coverage          & 268 instruments \\
    Date range        & 23~Feb~2025 -- 27~Feb~2026 \\
    Raw records       & 71,044 rows \\
    Clean records     & 70,029 rows \\
    Fields            & Symbol, date, open, high, low, close, volume \\
    Collection method & Python web scraping \\
    Benchmark index   & CSE All-Share Price Index (ASPI) \\
    Sectors           & 20 CSE standard sectors \\
    \bottomrule
  \end{tabular}
\end{table}

ASPI daily OHLC data were collected separately and merged on date to
provide the market benchmark return $R_{m}$ required by the market model.
Sector classifications were derived from CSE sector listings, mapping each
stock to one of 20 sectors.

\subsection{Data Cleaning}

\textbf{Suffix handling.} CSE tickers carry suffixes encoding share class.
Records with suffixes U0000 (units, $n=407$) and P0000 (preference shares,
$n=224$) were removed. Three symbols were subsequently removed explicitly:
Browns Beach Hotels (BBH) and Nation Lanka Finance (CSF) due to delisting
and data quality issues, and INME which held no observations within the
estimation window (Day~$-120$ to Day~$-6$). This yielded 70,029~rows
across voting (N0000, 65,279~rows) and non-voting (X0000, 4,807~rows)
shares.

\textbf{Sector mapping.} A manually verified sector mapping was applied to
the working dataset with 263 symbols.

\textbf{Regression qualification.} The market model was estimated only for
stocks with at least 60 clean trading observations in the estimation window
(Day~$-120$ to Day~$-6$), excluding 6 thinly traded stocks and leaving 257
stocks. Stocks with $\hat{\beta} \notin [-3,\,5]$ were subsequently removed
as outliers (20~stocks dropped), yielding a final sample of 237~stocks for
downstream analysis.

%%=============================================================================
\section{Methodology}
%%=============================================================================

We adopt the standard event-study framework of
MacKinlay~\cite{mackinlay1997}. This method estimates the return of each
stock absent the event and compares it with observed returns; the difference
is taken as the abnormal return, i.e.\ the portion of price movement
attributable directly to Cyclone Ditwah.

\subsection{Event Definition and Windows}

The event date ($t = 0$) is 28~November~2025 (cyclone landfall). Two
non-overlapping windows are defined:
\begin{itemize}
  \item \textbf{Estimation window:} Day~$-120$ to Day~$-6$
        (${\sim}115$ trading days). The gap between Day~$-6$ and Day~$-5$
        isolates the estimation period from potential anticipation effects.
  \item \textbf{Event window:} Day~$-5$ to Day~$+30$ (36 trading days).
        Spans pre-cyclone anticipation, the event shock, and post-disaster
        recovery.
\end{itemize}

\subsection{Market Model}

Individual stock returns are modelled as a linear function of market
returns over the estimation window:
\begin{equation}
  R_{it} = \alpha_i + \beta_i R_{mt} + \varepsilon_{it}
  \label{eq:market-model}
\end{equation}
where $R_{it}$ is the daily return of stock $i$ on day $t$, $R_{mt}$ is
the ASPI return, $\alpha_i$ is the stock-specific intercept, $\beta_i$ is
the market-sensitivity coefficient, and $\varepsilon_{it}$ is a zero-mean
error term.

\subsection{Abnormal Returns}

Abnormal returns in the event window are the residual between the actual
return and the counterfactual predicted by the fitted model:
\begin{equation}
  AR_{it} = R_{it} - \bigl(\hat{\alpha}_i + \hat{\beta}_i R_{mt}\bigr)
  \label{eq:ar}
\end{equation}
Cumulative abnormal returns aggregate over any sub-window $[t_1, t_2]$:
\begin{equation}
  CAR_i(t_1, t_2) = \sum_{t=t_1}^{t_2} AR_{it}
  \label{eq:car}
\end{equation}
We report CARs for four windows: pre-event $[-5,-1]$, landfall day
$[0,0]$, immediate aftermath $[0,+5]$, and the complete event window
$[-5,+30]$.

\subsection{Statistical Testing}
\label{sec:stat_testing}

\paragraph{Parametric cross-sectional $t$-test.}
To determine whether observed abnormal returns are statistically
distinguishable from zero, cross-sectional $t$-tests were applied at both
the market-wide and sector levels. For each window $[t_1,t_2]$, the
$CAR_i$ values across all stocks (or stocks within a sector) are treated as
a sample and tested against a null mean of zero:
\begin{equation}
  t = \frac{\overline{CAR}}{s / \sqrt{N}}
  \label{eq:tstat}
\end{equation}
where $\overline{CAR}$ is the cross-sectional mean, $s$ is the
cross-sectional standard deviation, and $N$ is the number of stocks. This
is a one-sample $t$-test against a population mean of zero (two-sided,
$\alpha = 0.05$).

\paragraph{Non-parametric Wilcoxon signed-rank test.}
Because financial return distributions in thin, emerging-market exchanges
routinely exhibit fat tails, skewness, and extreme outliers, relying
exclusively on the parametric $t$-test risks inflated Type~I error rates.
We therefore apply the Wilcoxon signed-rank test as a robust secondary
check on the full-window sector-level CAR distributions. The Wilcoxon test
evaluates whether the \emph{median} abnormal return is significantly
different from zero without assuming normality, making it insensitive to
the influence of a small number of extreme-performing equities. For each
sector with $N > 1$ stocks, the test statistic $W$ is computed from the
signed ranks of the $CAR_i$ values. Concordance between parametric and
non-parametric results strengthens confidence in the findings; divergence
signals that the sector's mean is driven by outliers rather than a uniform
industry-wide repricing.

\paragraph{Volatility and liquidity $t$-tests.}
To assess disaster-induced shifts in market microstructure, daily market
volatility was proxied using the High-Low price spread
($\text{Spread}_{t} = \text{High}_{t} - \text{Low}_{t}$), and market
liquidity was measured via total daily trading volume. Daily averages of
these metrics were computed across all qualifying stocks for two symmetric
25-day windows: a pre-event window (Day~$-30$ to Day~$-6$) and a
post-landfall window (Day~0 to Day~$+24$). Welch's independent-samples
$t$-test (unequal variances assumed) was used to compare the daily average
spread and daily average volume between the two windows.

%%=============================================================================
\section{Results}
%%=============================================================================

\subsection{ASPI Index Trajectory}

Figure~\ref{fig:aspi} plots the ASPI daily close from February~2025 to
March~2026. The index opened at approximately 16,700 in February~2025,
dipped to a local trough near 15,000 in late April~2025, then recovered
steadily to a pre-cyclone peak of approximately 23,600 in
mid-November~2025. Following landfall on 28~November~2025, the index
declined sharply from 22,844 (26~November) to a local trough of
approximately 21,431 on 5~December~2025—a peak-to-trough drop of
approximately 6.2\% over six trading days. The index subsequently
recovered, returning above 23,000 by January~2026 and reaching
approximately 24,000 by February~2026, suggesting the market largely
absorbed the cyclone shock within approximately 40 trading days.

\begin{figure}[h]
  \centering
  \includegraphics[width=\linewidth]{figure_aspi.pdf}
  \caption{ASPI daily close, Feb~2025--Mar~2026. Dashed vertical line
           marks landfall (28~Nov~2025). Shaded region indicates the
           acute decline period (Day~0 to Day~$+5$).}
  \label{fig:aspi}
\end{figure}

\subsection{Market Model Parameter Estimates}

Table~\ref{tab:params} summarises the distribution of $\hat{\alpha}$ and
$\hat{\beta}$ estimated from the estimation window. The mean $\hat{\beta}$
of 0.72 indicates that the average CSE stock moves less than the market,
consistent with relatively low liquidity and a high proportion of defensive
stocks. Figure~\ref{fig:beta} shows the pre-event beta distribution by
sector. Banks stocks exhibit the highest median beta, with the
interquartile range centred above 1.0. Transportation stocks show the
lowest median beta (approximately 0.2). Real Estate shows the widest
spread, reflecting high cross-stock heterogeneity within the sector.

\begin{table}[h]
  \centering
  \caption{Market model parameter estimates ($N = 237$ stocks)}
  \label{tab:params}
  \small
  \begin{tabular}{@{}lrrrr@{}}
    \toprule
    \textbf{Parameter} & \textbf{Mean} & \textbf{Std Dev} & \textbf{Min} & \textbf{Max} \\
    \midrule
    $\hat{\alpha}$ & $0.0009$ & $0.0031$ & $-0.0089$ & $0.0214$ \\
    $\hat{\beta}$  & $0.72$   & $0.48$   & $-0.31$   & $2.84$   \\
    \bottomrule
  \end{tabular}
\end{table}

\begin{figure}[h]
  \centering
  \includegraphics[width=\linewidth]{figure2_beta_by_sector.pdf}
  \caption{Pre-event beta distribution by sector (estimation window:
           Day~$-120$ to Day~$-6$).}
  \label{fig:beta}
\end{figure}

\subsection{Sector-Level Cumulative Abnormal Returns}

Table~\ref{tab:market_car} reports the market-wide mean CAR across all
237~qualifying stocks. The full-window mean CAR of $-4.42\%$ confirms a
broad market-level shock, with the largest single-period impact in the
immediate aftermath ($-1.71\%$ over Days~0 to~$+5$). The near-zero
landfall-day CAR$(0,0)$ of $-0.03\%$ indicates that the single-day shock
was negligible and that losses accumulated over subsequent days as damage
assessments emerged. The large cross-sectional standard deviation of
$11.39\%$ in the full window confirms high dispersion across individual
stocks, motivating the sector-level disaggregation below.

\begin{table}[h]
  \centering
  \caption{Market-level mean CAR ($n=237$ stocks except Day~0 where
           $n=233$ due to non-trading days)}
  \label{tab:market_car}
  \small
  \begin{tabular}{@{}lrr@{}}
    \toprule
    \textbf{Window} & \textbf{Mean CAR (\%)} & \textbf{Std Dev (\%)} \\
    \midrule
    Pre-event $[-5,-1]$      & $-0.78$ & $4.56$  \\
    Landfall day $[0,0]$     & $-0.03$ & $2.65$  \\
    Aftermath $[0,+5]$       & $-1.71$ & $5.95$  \\
    Full window $[-5,+30]$   & $-4.42$ & $11.39$ \\
    \bottomrule
  \end{tabular}
\end{table}

Figure~\ref{fig:car} plots the mean CAR trajectory by sector across the
full event window, anchored at Day~0. Table~\ref{tab:car} reports point
estimates for all 20~CSE industry groups.

\begin{figure}[h]
  \centering
  \includegraphics[width=\linewidth]{figure1_sector_CAR.pdf}
  \caption{Sector-level cumulative abnormal returns anchored at Day~0
           (landfall, 28~Nov~2025). Pre-event values show reverse
           accumulation back from Day~0; post-event values show forward
           accumulation from Day~0. Shaded region = anticipation window
           Day~$-5$ to Day~0.}
  \label{fig:car}
\end{figure}

\begin{table}[!htb]
  \centering
  \caption{Sector-level CAR estimates (\%)}
  \label{tab:car}
  \footnotesize
  \resizebox{\columnwidth}{!}{%
  \begin{tabular}{@{}lrrrr@{}}
    \toprule
    \textbf{Sector} & \textbf{CAR($-5,-1$)} & \textbf{CAR($0,0$)} &
    \textbf{CAR($0,+5$)} & \textbf{CAR($-5,+30$)} \\
    \midrule
    Transportation         & $-4.53$ & $-1.21$ & $-4.45$ & $-10.03$ \\
    Real Estate            & $-2.36$ & $-0.18$ & $-2.57$ & $-9.23$  \\
    Software \& Services   & $-1.26$ & $ 0.06$ & $-2.26$ & $-7.41$  \\
    Utilities              & $ 0.77$ & $-0.86$ & $-1.64$ & $-6.69$  \\
    Diversified Financials & $-0.96$ & $-0.34$ & $-2.44$ & $-6.69$  \\
    Insurance              & $-0.97$ & $-0.70$ & $-4.71$ & $-6.36$  \\
    Banks                  & $-1.91$ & $ 0.09$ & $-1.17$ & $-6.07$  \\
    Telecom Services       & $-3.53$ & $ 0.36$ & $-1.98$ & $-5.34$  \\
    Food \& Beverage       & $-0.87$ & $ 0.19$ & $-1.76$ & $-4.18$  \\
    Consumer Services      & $ 0.90$ & $-0.33$ & $-3.83$ & $-3.97$  \\
    Capital Goods          & $-1.70$ & $ 0.12$ & $-0.22$ & $-2.97$  \\
    Retailing              & $-0.92$ & $ 2.25$ & $-2.13$ & $-2.50$  \\
    Consumer Durables      & $ 0.24$ & $-0.25$ & $ 0.65$ & $-2.28$  \\
    Energy                 & $-0.70$ & $ 0.05$ & $-0.69$ & $-2.09$  \\
    Materials              & $-1.55$ & $-0.43$ & $ 0.43$ & $-1.75$  \\
    Household Products     & $-0.79$ & $-0.08$ & $-1.66$ & $-1.22$  \\
    Healthcare Equipment   & $-0.89$ & $ 0.58$ & $ 1.89$ & $-0.66$  \\
    Food Retailing         & $-0.02$ & $-0.34$ & $ 1.54$ & $-0.04$  \\
    Commercial Services    & $ 4.66$ & $-0.73$ & $ 1.87$ & $ 0.27$  \\
    Automobiles            & $ 1.42$ & $ 0.16$ & $-0.66$ & $ 2.43$  \\
    \bottomrule
  \end{tabular}}
\end{table}

\textbf{Most affected sectors.}
Transportation recorded the largest full-window CAR at $-10.03\%$, with
its trajectory reaching a trough of approximately $-12\%$ around Day~$+26$
before a partial recovery. Real Estate followed at $-9.23\%$, consistent
with the cyclone's documented damage to 247~km of roads, 35~bridges, and
coastal property. Insurance recorded the steepest immediate aftermath
CAR$(0,+5)$ at $-4.71\%$, reflecting anticipated claims exposure.
Software \& Services ($-7.41\%$) and Utilities ($-6.69\%$) also sustained
substantial full-window losses, likely reflecting infrastructure disruption
and power outage effects.

\textbf{Anticipation effects.}
In the pre-event window, Transportation shows a CAR$(-5,-1)$ of $-4.53\%$
and Telecommunication Services $-3.53\%$—the two largest pre-event
negative CARs across all 20~sectors—confirming that abnormal losses were
accumulating before official landfall as investors reacted to cyclone
warnings.

\textbf{Resilient sectors.}
Materials remained near zero throughout ($-1.75\%$). Commercial Services
($+0.27\%$) and Automobiles ($+2.43\%$) posted near-zero or positive
full-window CARs, possibly reflecting reconstruction-driven demand. On
landfall day alone, Retailing recorded an anomalous $+2.25\%$, the only
sector with a substantial positive single-day abnormal return, possibly
reflecting panic-buying ahead of the storm.

\subsection{Sector-Level Significance: Parametric and Non-Parametric Tests}

Table~\ref{tab:significance} reports both the cross-sectional $t$-test and
the Wilcoxon signed-rank test for the full event window CAR$(-5,+30)$
across all 20~sectors. Figure~\ref{fig:car_boxplot} shows the within-sector
CAR distributions that underpin the significance results.

%% FIX: boxplot figure inserted here, where it is first referenced
\begin{figure}[h]
  \centering
  \includegraphics[width=\linewidth]{sector_car_boxplot.png}
  \caption{Distribution of full-window CAR$(-5,+30)$ by sector
           (box = IQR, whiskers = 1.5$\times$IQR, circles = outliers).
           Sectors ordered by median CAR. The x-axis extends to
           $\pm400$--$600\%$ due to extreme individual-stock outliers in
           Diversified Financials, Food Retailing, Insurance, and
           Materials. Despite Transportation recording the largest
           negative mean CAR ($-10.03\%$), its box is tight and centred
           near zero, confirming that the mean is driven by one or two
           extreme equities rather than a uniform crash. Real Estate's
           box sits clearly left of zero across most of its distribution,
           consistent with its statistically significant systemic
           repricing.}
  \label{fig:car_boxplot}
\end{figure}

\begin{table*}[!htb]
  \centering
  \caption{Sector-level statistical significance: cross-sectional $t$-test
           and Wilcoxon signed-rank test, full event window CAR$(-5,+30)$.
           Sectors ordered by mean CAR. $^{\dagger}$Sectors with $N < 5$
           (Energy $N=2$; Software \& Services $N=1$; Automobiles $N=1$;
           Telecommunication Services $N=2$; Food Retailing $N=3$) are
           included for completeness but should be interpreted with caution;
           test statistics are not computed for $N=1$.}
  \label{tab:significance}
  \footnotesize
  \begin{tabular}{@{}lrrrrrrr@{}}
    \toprule
    \multirow{2}{*}{\textbf{Sector}} &
    \multirow{2}{*}{\textbf{N}} &
    \multicolumn{3}{c}{\textbf{$t$-test (mean CAR)}} &
    \multicolumn{3}{c}{\textbf{Wilcoxon (median CAR)}} \\
    \cmidrule(lr){3-5}\cmidrule(lr){6-8}
    & & \textbf{Mean (\%)} & \textbf{$t$} & \textbf{$p$} &
        \textbf{Median (\%)} & \textbf{$W$} & \textbf{$p$} \\
    \midrule
    Transportation$^{\dagger}$       & 3  & $-10.03$ & $-1.16$ & $0.366$ & $-17.33$ & $1$  & $0.500$ \\
    Real Estate                      & 17 & $ -9.23$ & $-2.55$ & $\mathbf{0.021}$ & $-5.43$ & $27$ & $\mathbf{0.017}$ \\
    Software \& Services$^{\dagger}$ & 1  & $ -7.41$ & ---     & ---     & $-7.41$ & ---  & ---     \\
    Utilities                        & 7  & $ -6.69$ & $-1.29$ & $0.244$ & $-0.18$ & $14$ & $1.000$ \\
    Diversified Financials           & 38 & $ -6.69$ & $-0.82$ & $0.420$ & $-5.79$ & $211$ & $\mathbf{0.034}$ \\
    Insurance                        & 12 & $ -6.36$ & $-1.20$ & $0.257$ & $-8.57$ & $7$  & $\mathbf{0.009}$ \\
    Banks                            & 13 & $ -6.07$ & $-1.25$ & $0.236$ & $-1.49$ & $45$ & $1.000$ \\
    Telecom Services$^{\dagger}$     & 2  & $ -5.34$ & $-1.71$ & $0.337$ & $-5.34$ & $0$  & $0.500$ \\
    Food \& Beverage                 & 43 & $ -4.18$ & $-0.50$ & $0.619$ & $-2.56$ & $326$ & $0.077$ \\
    Consumer Services                & 32 & $ -3.97$ & $-0.14$ & $0.887$ & $-1.10$ & $203$ & $0.262$ \\
    Capital Goods                    & 29 & $ -2.97$ & $-0.27$ & $0.787$ & $-3.03$ & $202$ & $0.749$ \\
    Retailing                        & 14 & $ -2.50$ & $-1.12$ & $0.282$ & $-2.55$ & $40$ & $0.463$ \\
    Consumer Durables                & 8  & $ -2.28$ & $-0.62$ & $0.555$ & $ 1.89$ & $16$ & $0.844$ \\
    Energy$^{\dagger}$               & 2  & $ -2.09$ & $-1.04$ & $0.487$ & $-52.38$ & $0$ & $0.500$ \\
    Materials                        & 21 & $ -1.75$ & $-0.72$ & $0.483$ & $-2.84$ & $77$ & $0.296$ \\
    Household Products               & 2  & $ -1.22$ & $-1.00$ & $0.500$ & $-0.61$ & $0$  & $1.000$ \\
    Healthcare Equipment             & 8  & $ -0.66$ & $-0.11$ & $0.914$ & $ 1.89$ & $15$ & $0.742$ \\
    Food Retailing$^{\dagger}$       & 3  & $ -0.04$ & $-1.00$ & $0.423$ & $ 4.58$ & $2$  & $0.750$ \\
    Commercial Services              & 7  & $ +0.27$ & $-0.44$ & $0.675$ & $-7.62$ & $8$  & $0.375$ \\
    Automobiles$^{\dagger}$          & 1  & $ +2.43$ & ---     & ---     & $ 2.43$ & ---  & ---     \\
    \midrule
    \textbf{All Stocks (Market)} & \textbf{237} & $\mathbf{-4.42}$ &
    $\mathbf{-6.14}$ & $\mathbf{<0.001}$ & --- & --- & --- \\
    \bottomrule
  \end{tabular}
\end{table*}

\textbf{Parametric results.}
Cross-sectional $t$-tests confirm a highly significant market-wide shock
($t = -6.14$, $p < 0.001$, $N = 237$). At the sector level, only Real
Estate achieves significance at the 5\% level ($t = -2.55$, $p = 0.021$,
$N = 17$). Transportation, despite the largest mean CAR of $-10.03\%$,
does not reach significance ($t = -1.16$, $p = 0.366$, $N = 3$),
reflecting both the small sample and the tight, near-zero box visible in
Figure~\ref{fig:car_boxplot} that confirms the mean is outlier-driven.
Sectors such as Diversified Financials and Consumer Services also show
non-significant results despite moderate negative means, driven by large
cross-sectional dispersion among their many constituent stocks.

\textbf{Non-parametric corroboration.}
The Wilcoxon signed-rank test reveals three sectors with statistically
significant median CARs at the 5\% level. Insurance is the most
significant ($W = 7$, $p = 0.009$), with a median CAR of $-8.57\%$,
reflecting broad-based claims exposure across all 12 constituent stocks.
Real Estate ($W = 27$, $p = 0.017$, median $-5.43\%$) and Diversified
Financials ($W = 211$, $p = 0.034$, median $-5.79\%$) are also
significant. Notably, Diversified Financials passes the Wilcoxon test
despite a non-significant $t$-test ($p = 0.420$), indicating that its
mean CAR is suppressed by a large positive outlier (visible in
Figure~\ref{fig:car_boxplot}) while the majority of its 38 stocks
experienced genuine negative abnormal returns. For Transportation, the
Wilcoxon test fails to reject the null ($W = 1$, $p = 0.500$),
consistent with the mean being outlier-driven. The divergence between
the $t$-test and Wilcoxon results for Diversified Financials is
precisely the kind of finding the non-parametric test is designed to
surface: when a sector's mean is distorted by a single extreme equity,
the Wilcoxon median test provides a more reliable picture of the
typical stock-level experience.

\subsection{Market Volatility and Liquidity Dynamics}

Beyond asset pricing, Cyclone Ditwah induced distinct and temporally
nuanced shifts in market microstructure. Table~\ref{tab:vol_liq}
summarises the aggregate pre- and post-event comparison, while
Figures~\ref{fig:vol_ts} and~\ref{fig:liq_ts} reveal the day-by-day
dynamics that the aggregate statistics partially conceal.

\begin{table}[h]
  \centering
  \caption{Pre- vs.\ post-event volatility and liquidity
           (Welch's independent-samples $t$-test;
           pre: Day~$-30$ to Day~$-6$;
           post: Day~0 to Day~$+24$; $n=25$ daily observations per window)}
  \label{tab:vol_liq}
  \small
  \begin{tabular}{@{}lrrrr@{}}
    \toprule
    \textbf{Metric} & \textbf{Pre-event} & \textbf{Post-event} &
    \textbf{$t$} & \textbf{$p$} \\
    \midrule
    Mean H-L spread (LKR)      & 10.77   & 8.87    & $1.84$ & $0.076$   \\
    Mean daily vol.\ (shares)  & 565{,}798 & 366{,}999 & $3.76$ & $<0.001$ \\
    \bottomrule
  \end{tabular}
\end{table}

%% FIX: filenames corrected — volatility (spread) figure uses
%% market_volatility_lineplot.png; liquidity (volume) uses
%% market_liquidity_lineplot.png
\begin{figure}[h]
  \centering
  \includegraphics[width=\linewidth]{market_volatility_lineplot.png}
  \caption{Average daily High-Low price spread (market volatility proxy),
           Day~$-30$ to Day~$+30$. The spread reaches its window-high of
           approximately 17.8~LKR on Day~0 itself, then collapses and
           trends downward through Day~$+30$. The post-event aggregate
           mean (8.87~LKR) is lower than the pre-event mean (10.77~LKR)
           because the Day~0 spike is immediately followed by a sustained
           compression, not because the disaster suppressed intraday price
           movement uniformly.}
  \label{fig:vol_ts}
\end{figure}

\begin{figure}[h]
  \centering
  \includegraphics[width=\linewidth]{market_liquidity_lineplot.png}
  \caption{Total daily trading volume (market liquidity proxy),
           Day~$-30$ to Day~$+30$ (y-axis in units of $10^8$ shares).
           Volume begins declining from approximately Day~$-10$, reaches
           its window minimum on Day~0 (${\sim}60$~million shares), briefly
           rebounds on Days~$+2$ to~$+4$, then contracts again and remains
           depressed through Day~$+30$, never recovering to pre-event
           levels within the observation window.}
  \label{fig:liq_ts}
\end{figure}

\textbf{Volatility: a Day-0 spike followed by sustained compression.}
Figure~\ref{fig:vol_ts} reveals a pattern the aggregate $t$-test statistic
alone does not capture. The High-Low spread reaches approximately
17.8~LKR on Day~0—the highest value in the entire Day~$-30$ to Day~$+30$
window—indicating a sharp burst of intraday price uncertainty at the moment
of landfall. This spike is immediately followed by a sustained downward
trend: by Day~$+10$ the spread has fallen below 11~LKR, and by Day~$+25$
it stabilises near 5--6~LKR, well below the pre-event baseline. As a
result, the post-event window aggregate mean (8.87~LKR) is actually
\emph{lower} than the pre-event mean (10.77~LKR), and the Welch $t$-test
of this aggregate difference is not significant at the 5\% level
($t = 1.84$, $p = 0.076$). The correct interpretation is a two-phase volatility
response: an acute landfall-day uncertainty spike, followed by progressive
compression as thinning liquidity reduced the range of transactions
executed each day.

\textbf{Liquidity: pre-event deterioration and a persistent post-event freeze.}
Figure~\ref{fig:liq_ts} reveals that the volume decline began approximately
10 trading days before landfall—consistent with the anticipation effects
observed in sector-level CARs—and reached a window minimum of approximately
60 million shares on Day~0. This pre-landfall deterioration is itself
informative: investors began reducing market participation as cyclone
warnings intensified, creating a liquidity drought that preceded the
physical disaster. Volume briefly rebounded on Days~$+2$ to~$+4$ (to
approximately 225 million shares), likely reflecting opportunistic or
rebalancing trades, before collapsing again and remaining well below
pre-event levels through Day~$+30$. The aggregate decline of approximately
35\% is highly significant ($t = 3.76$, $p < 0.001$).

\textbf{Interpretation: a two-phase microstructure response.}
Taken together, the figures point to a two-phase response. In the first
phase (approximately Day~$-10$ to Day~$+1$), anticipation and landfall
drove simultaneous liquidity withdrawal and a brief but sharp intraday
volatility spike. In the second phase (Day~$+2$ onward), the market
entered a \emph{liquidity freeze}: volume remained severely depressed and
intraday price ranges compressed as uncertainty suppressed both demand and
supply of shares. This pattern contrasts with a panic-selling regime, which
would exhibit elevated volume \emph{and} sustained volatility. The policy
implication is phase-specific: circuit breakers are most relevant at the
Day~0 spike, while liquidity-restoration measures—market-maker incentives,
temporary transaction cost reductions—are needed to address the
longer-duration freeze.

%%=============================================================================
\section{Discussion}
%%=============================================================================

\textbf{Market-wide shock.}
The full-window market mean CAR of $-4.42\%$ across 237~qualifying stocks,
confirmed by a highly significant $t$-statistic ($t = -6.14$, $p < 0.001$),
establishes that Cyclone Ditwah induced a meaningful aggregate shock on the
CSE. The near-zero landfall-day CAR$(0,0)$ of $-0.03\%$ is striking: the
single-day impact was negligible market-wide, and the bulk of losses
accumulated over subsequent days as damage assessments emerged. This
pattern of gradual repricing is consistent with the semi-strong form of the
Efficient Market Hypothesis, in which successively released public
information is incorporated incrementally into prices.

\textbf{Anticipation effects.}
Transportation's pre-event CAR$(-5,-1)$ of $-4.53\%$ and Telecom Services'
$-3.53\%$ are the two largest pre-event negative CARs across all
20~sectors, suggesting investors began pricing in disaster risk as cyclone
warnings intensified before official landfall. These anticipation patterns
are consistent with semi-strong form market efficiency
\cite{mackinlay1997}. Notably, the volume time series
(Figure~\ref{fig:liq_ts}) corroborates this: liquidity began deteriorating
from approximately Day~$-10$, simultaneous with the onset of pre-event
abnormal returns in vulnerable sectors.

\textbf{Sector heterogeneity and parametric/non-parametric divergence.}
Transportation ($-10.03\%$) and Real Estate ($-9.23\%$) bore the greatest
cumulative mean losses, consistent with the cyclone's documented damage to
247~km of roads, 35~bridges, and coastal property \cite{who2025, ocha2025}.
However, the Wilcoxon test identifies three sectors with statistically
significant median repricing: Insurance ($p = 0.009$), Real Estate
($p = 0.017$), and Diversified Financials ($p = 0.034$). Transportation,
despite its large negative mean, fails both tests—its box in
Figure~\ref{fig:car_boxplot} is tight and centred near zero, confirming
the mean is outlier-driven across only three stocks. The most instructive
divergence is Diversified Financials: its $t$-test is non-significant
($p = 0.420$) because a single large positive outlier inflates the mean,
yet the Wilcoxon test reveals the median stock in the sector did experience
significant negative abnormal returns ($p = 0.034$). Insurance shows the
opposite pattern—a non-significant $t$-test ($p = 0.257$) but the
strongest Wilcoxon result ($p = 0.009$), driven by consistently negative
returns across its 12 stocks rather than a large mean. These divergences
illustrate why the non-parametric test is essential in thin, heterogeneous
markets: the mean CAR is an unreliable guide to whether a typical investor
in a sector was affected.

Insurance recorded the steepest immediate aftermath CAR$(0,+5)$ at
$-4.71\%$, reflecting anticipated claims exposure. Software \& Services
($-7.41\%$) and Utilities ($-6.69\%$) also sustained substantial losses,
likely reflecting infrastructure disruption. Materials ($-1.75\%$) remained
near zero throughout, and Commercial Services ($+0.27\%$) and Automobiles
($+2.43\%$) posted positive full-window CARs, possibly reflecting
reconstruction-driven demand.

\textbf{Index versus sector recovery divergence.}
The ASPI returned above 23,000 by January~2026, suggesting broad market
recovery within approximately 40~trading days. However, sector-level CARs
at Day~$+30$ remain substantially negative for Transportation, Real Estate,
and Software \& Services, indicating that index-level recovery masked
persistent sector-specific damage. Headline index recovery does not imply
uniform stock-level recovery, and sector-disaggregated monitoring is
necessary to capture residual disaster exposure.

\textbf{Market microstructure and the two-phase response.}
The microstructure evidence enriches the asset pricing results with a
temporal dimension. The High-Low spread time series reveals a sharp
intraday volatility spike on Day~0 (${\sim}17.8$~LKR, the window
maximum), immediately followed by a sustained compression to well below
pre-event levels. Simultaneously, the volume time series shows that the
liquidity deterioration began approximately 10 days before landfall, hit
its minimum at Day~0, and never recovered within the observation window.
Together these patterns describe a two-phase response: a brief acute shock
at landfall followed by a persistent liquidity freeze. The policy
implication is phase-specific: circuit breakers are most relevant at the
Day~0 spike, whereas longer-duration liquidity-restoration tools are
required to address the extended freeze that follows.

%%=============================================================================
\section{Limitations and Future Work}
%%=============================================================================

\paragraph{Single-factor market model.}
The market model regresses individual stock returns on the ASPI alone. A
multi-factor specification incorporating size, value, and momentum factors
(as in the Fama-French framework) would provide a richer counterfactual and
may reduce omitted-variable bias. Future work will also incorporate
macroeconomic control variables from the Central Bank of Sri Lanka and IMF
country reports \cite{imf2025}.

\paragraph{Volatility proxy.}
The High-Low daily price spread is a widely used but imperfect proxy for
realised volatility. While it successfully captured the Day~0 intraday
uncertainty spike (${\sim}17.8$~LKR) and subsequent compression, it does
not account for intraday price paths, bid-ask bounce, or overnight gaps.
Access to tick-by-tick data would permit the use of realised variance
estimators and would better decompose the Day~0 spike into its fundamental
versus noise components.

\paragraph{Small-sector reliability.}
Five sectors contain fewer than five stocks (Energy $N=2$, Automobiles
$N=1$, Software \& Services $N=1$, Telecommunication Services $N=2$, Food
Retailing $N=3$). Point estimates are reported for completeness but carry
very wide confidence intervals; their statistical tests are either
underpowered or not computable. Future work should aggregate these into
broader categories or expand coverage to OTC-traded instruments.

\paragraph{Sample period and confounding events.}
The estimation window spans a period of ongoing IMF programme negotiations
and macroeconomic uncertainty \cite{imf2025}. Firm-specific confounders
during the event window—earnings releases, rights issues, management
changes—cannot be individually controlled without a comprehensive corporate
actions database, which future CSE event studies should incorporate.

\paragraph{Generalisation.}
This study examines a single disaster event on a single frontier exchange.
Caution is warranted in generalising to other disaster types or markets.
Cross-event and cross-market replication studies would establish whether the
liquidity-freeze pattern and sector heterogeneity documented here are robust
features of emerging-market disaster responses.

%%=============================================================================
\begin{thebibliography}{9}

\bibitem{worldbank2025}
World Bank.
\newblock Damage from Cyclone Ditwah in Sri Lanka Estimated at \$4.1 Billion.
\newblock \textit{World Bank Press Release}, December~22, 2025.

\bibitem{who2025}
World Health Organization.
\newblock Sri Lanka Floods and Landslides -- Cyclonic Storm Ditwah
  November~2025.
\newblock \textit{WHO South-East Asia Region}, December~2, 2025.

\bibitem{ocha2025}
OCHA.
\newblock Sri Lanka: Cyclone Ditwah -- Situation Report~3.
\newblock \textit{UN Office for the Coordination of Humanitarian Affairs},
  December~23, 2025.

\bibitem{brecorder2025}
Business Recorder.
\newblock Cyclone-ravaged Sri Lanka's shares log worst week in over three
  years.
\newblock 2025.

\bibitem{imf2025}
International Monetary Fund.
\newblock IMF Country Report No.\ 25/339: Sri Lanka.
\newblock 2025.

\bibitem{mackinlay1997}
A.~C. MacKinlay.
\newblock Event studies in economics and finance.
\newblock \textit{Journal of Economic Literature}, 35(1):13--39, 1997.

\end{thebibliography}

\end{document}